\documentclass[main.tex]{subfiles}
\begin{document}
For the finite-dimensional complex Lie algebras $\mathfrak{sp}(2n)$ and $\mathfrak{so}(n)$, all finite-dimensional representations form a semisimple category according to Hermann Weyl's semisimplicity theorem. However, in the case of the classical infinite-dimensional locally finite Lie algebras  $\mathfrak{sp}_\infty$ and $\mathfrak{so}_\infty$, there are no non-trivial finite-dimensional representations. Natural analogues of finite-dimensional representations are the tensor representations and, as shown in \cite{penkov2007tensorrepresentationsclassicallocally}, these do not form semisimple categories.\\ 

Indecomposable injective objects, as well as injective resolutions of simple objects in the respective categories of tensor $\mathfrak{sp}_\infty- $ and $ \mathfrak{so}_\infty-$modules have been studied in detail (\cite{penkov2007tensorrepresentationsclassicallocally}, \cite{dancohen2012koszulcategoryrepresentationsfinitary}, \cite{SAM_2015}). The published explicit formulas for the layers of the socle filtrations of injective objects and for terms of injective resolutions involve Littlewood-Richardson coefficients. These coefficients are computed algorithmically and require several steps for checking whether or not they are nonzero.
Hence the socle layers of injective tensor modules, as well as the terms of injective resolutions, are difficult to study without the usage of a computer program with an implemented Littlewood-Richardson coefficient calculator.\\

In this Bachelor's thesis, we are simplifying some aspects of the already known theory, by producing direct descriptions of modules appearing in the top layers of the socle filtrations of injective $\mathfrak{sp}_\infty-,\mathfrak{so}_\infty-$modules, and in the last terms of minimal injective resolutions of $\mathfrak{sp}_\infty-,\mathfrak{so}_\infty-$modules. Additionally, we provide explicit formulas for the length of the socle filtrations and the length of minimal injective resolutions of such modules. These formulas do not depend on Littlewood-Richardson coefficients and are written directly in terms of the partition $\lambda$ that determines the respective $\mathfrak{sp}_\infty-$, $\mathfrak{so}_\infty-$module $\Gamma_\lambda$ (as defined in section 2).  \\

On an organizational note, this thesis starts with introducing the necessary background in combinatorics and algebra, using $\cite{fulton1991representation}$ and $\cite{penkov2007tensorrepresentationsclassicallocally}$ as main references for the theory. Section 3 includes a discussion of socle filtrations, using the theorem describing the socle layers of an injective tensor
$\mathfrak{sp}_\infty-$module (\cite{penkov2007tensorrepresentationsclassicallocally}) as a key starting point. We provide a more explicit version of this theorem by finding the length of the socle filtration, and then describe some direct summands of the last socle layer, that are guaranteed to appear. Section 4 presents an algorithm for computing injective resolutions, a combinatorial result involving column-quasi-symmetric partitions, and some consequences of the description of the terms of minimal injective resolutions of simple tensor $\mathfrak{sp}_\infty-$modules given by \cite{SAM_2015}. In Section 5, we showcase the analogous results for the Lie algebra $\mathfrak{so}_\infty.$ We complete the thesis by a brief conclusion and a possible future research question.\\

As an appendix we present computer-generated examples of minimal injective resolutions of simple tensor $\mathfrak{sp}_\infty-$modules.
\bigskip                                                                     
\end{document}


